\section{スケール反転：個体から社会へ}
\label{sec:multiscale}

\subsection{機能と支配的コストの比較}
前節までのモデルは、個体の認知系のレベルで解釈している。
その解釈では、状態を再訪することと、再利用可能な離散的な手がかりで連続配置を表現すること、すなわち再帰と圧縮を支配的なコストとして提案する。
こうした構成上の仮定のもとで、言語トークンに内部制御の役割を割り当てる。

この役割を相互作用するエージェント間の同期・協調と比較することで、社会的拡張を提案する。
トークンは共有軌道の安定化を助け、他のエージェントの反応についての不確実性を減らすと仮定する。
表は提案する支配的コストとトークンの機能を対比するものであり、個体と集合のダイナミクスが同一であることを確立しない。

\begin{center}
\begin{tabular}{lll}
\toprule
\textbf{スケール} & \textbf{支配的コスト} & \textbf{トークンの機能} \\
\midrule
個体 & 再帰、圧縮 & 内部制御単位 \\
社会 & 同期、協調 & 共有地形のアンカー \\
\bottomrule
\end{tabular}
\end{center}

この比較は、異なる有効制約のもとで類似した再アドレス機能を提案する。
スケールをまたいでトークンの性質が不変であることを確立するものではない。
集合的な集約、共通の状態空間、型を明示した翻訳は未指定であり、同じ記法の再利用から共有場の測定が得られるわけではない。

\subsection{社会的な概念地形}
社会的仮説では、事前の地形が十分に類似しているとき、共有された言語入力への反復的な接触を、相関した井戸の深化と比較する。
共有概念地形は、この提案された組織化に対するモデルレベルの名称である。

この描像では、共有地形は一人のエージェント内に保存されるのではなく、エージェント間に分布し、継続的な言語的相互作用を通じて維持される。
履歴によって形成された集合的幾何は、これらの相互作用の解釈であり、測定された集合場ではない。

\subsection{規範、制度、パラダイム}
深く沈殿した社会的領域を、規範、制度、科学的パラダイムと構造的に比較する。
このアトラクタ類比では、集合的な$\Phi$を用いて持続を表現するが、制度と形式的アトラクタの同一性を主張しない。

別の仮説として、深さを安定性と逸脱コストに、協調した摂動を制度変化に関連づける。
これは中央による強制を唯一の説明とせずに持続を考える方法を提供するものであり、強制が不要であることの証拠ではない。
パラダイム転換の歴史的困難さは比較対象の候補であって、井戸の深さによって説明された結果ではない。
独立した制度の基準と深さの推定は未解決である。
安定性を深さの定義と、その説明上の役割を支える独立した証拠の両方に用いることはできない。

\subsection{地形維持としての知性}
この社会的スケールの記述内では、提案する分散的過程に対して局所的な機能上の呼称を用いる。
\begin{quote}
\itshape
ここでいう社会的知性とは、言語を通じた概念地形の分散的な維持、同期、再形成の呼称である。
\end{quote}

この命名は、知性の確立された一般的定義や生成ではなく、個体の知性を排するものでもない。
repository内のIntelligence諸理論との関係は、概念的な引き継ぎであって、下流の研究から検証を移し入れるものではない。

個々のエージェントは、この過程の一時的な参加者としてモデル化する。
地形を通過し、局所的な変形を導入し、外在化された言語ログを通じて安定化した構造を後続のエージェントへ伝える。

厳密な正値性$I>0$は、提案された社会的地形について類比的な非閉包の読みを動機づけるものであり、集合的ダイナミクスの導出ではない。
さらなる分化と革新は、残留非対称性を再活性化する新たな摂動や、より広い観測窓のもとで提案する可能性であり、正値性だけで保証される結果ではない。
